\section{Introduction}

Sparse mixture-of-experts (MoE) layers increase parameter capacity by routing
each token to a small subset of experts \citep{shazeer2017outrageously,
fedus2022switch}. In a linear top-1 router, a token hidden state is assigned to
the expert whose router row has the largest inner product with it. Therefore,
when we ask whether uniformly occurring features can be routed to uniformly used
experts, the immediate mathematical object is not the language task, but the
geometry seen by this inner product.

The controlled setting is deliberately simpler than real language modeling.
Consider many independent feature families with equal frequency and no semantic
relation across families. Even in this simplified regime, several notions of
``uniform'' must be separated. Feature labels can be frequency-balanced while
the hidden states are not uniformly distributed in high-dimensional space. The
hidden covariance can have no dominant eigendirection while the direction
distribution is still not spherical. Expert usage counts can be forced to be
exactly equal by a capacity rule, but this does not explain why top-1 routing
would balance naturally. Conversely, router rows can align with expert singular
directions while still producing imbalanced usage when those directions cluster.

We restrict the paper to the static controlled case. We do not claim that a
short-context next-token-prediction model naturally learns uniform hidden
states; such states may be strongly anisotropic because of local prediction
signals and residual common components. Instead, we ask: under what
representation-geometric assumptions is uniform expert usage a consequence of
unconstrained top-1 routing itself?

\paragraph{Main message.}
Uniform top-1 expert usage is a two-sided sufficient-condition result. On the
token side, the route-relevant feature directions must be close to a spherical
uniform distribution. On the expert side, the router rows or expert-derived
prototypes must cut that sphere into equal-measure cells. If either side fails,
usage can be imbalanced even when the other side is well behaved.

\paragraph{Contributions.}
\begin{enumerate}
\item We define the controlled uniform feature routing problem and separate
feature-frequency uniformity, hidden spectral uniformity, directional
uniformity, and expert-usage uniformity.
\item We prove a sufficient condition for natural top-1 usage balance, decompose
finite-token imbalance, and show why covariance isotropy alone is not a routing
certificate.
\item We validate the theory with synthetic experiments, treating expert
spectral prototypes as center construction and balanced assignment as a
non-geometric control.
\end{enumerate}
